\begin{examplebox}
	\textbf{\textcolor{CExample}{例 1（基础）}}
	\textbf{\textcolor{CExample}{题目：}} 已知函数 $f(x)=x^2-2x+2$ 在区间 $[0,3]$ 上，若 $\forall x\in[0,3],\ f(x)\geq a$ 恒成立，求实数 $a$ 的取值范围。\textbf{\textcolor{CExample}{解题步骤：}}
	\begin{description}[style=nextline, leftmargin=2.8em, labelsep=0.5em, itemsep=0.45em]
		\item[\textbf{第一步（最小值）：}]
			配方法 $f(x)=(x-1)^2+1$，最小值 $m=1$。
			\par\textbf{理由：} 二次函数开口向上，对称轴在区间内。
		\item[\textbf{第二步（全称转化）：}]
			由结论 $\forall f\geq a \iff m\geq a$，得 $1\geq a$，即 $a\leq 1$。
			\par\textbf{理由：} 直接套用全称量词转化为最值。
		\item[\textbf{第三步（答案）：}]
			$a\in(-\infty,1]$。
			\par\textbf{理由：} 解不等式。
	\end{description}
	\textbf{\textcolor{CExample}{关键结论：}}
	\textbf{$a\leq 1$}

	\medskip
	\textbf{\textcolor{CExample}{例 2（稍变形）}}
	\textbf{\textcolor{CExample}{题目：}} 已知函数 $f(x)=x^2-2x+2$ 在 $[0,3]$ 上，若 $\exists x\in[0,3],\ f(x)\leq a$，求 $a$ 的取值范围。\textbf{\textcolor{CExample}{解题步骤：}}
	\begin{description}[style=nextline, leftmargin=2.8em, labelsep=0.5em, itemsep=0.45em]
		\item[\textbf{第一步（最小值同）：}]
			$m=1$。
			\par\textbf{理由：} 同例1。
		\item[\textbf{第二步（存在转化）：}]
			$\exists f\leq a \iff m\leq a$，得 $a\geq 1$。
			\par\textbf{理由：} 存在量词转化为最小值不大于参数。
		\item[\textbf{第三步（答案）：}]
			$a\in[1,+\infty)$。
			\par\textbf{理由：} 解 $a\geq 1$。
	\end{description}
	\textbf{\textcolor{CExample}{关键结论：}}
	\textbf{$a\geq 1$}

	\medskip
	\textbf{\textcolor{CExample}{例 3（综合应用）}}
	\textbf{\textcolor{CExample}{题目：}} 已知函数 $f(x)=x^2+ax+1$ 在 $[-1,1]$ 上，若 $\forall x\in[-1,1],\ f(x)\leq 5$ 恒成立，求 $a$ 的范围。\textbf{\textcolor{CExample}{解题步骤：}}
	\begin{description}[style=nextline, leftmargin=2.8em, labelsep=0.5em, itemsep=0.45em]
		\item[\textbf{第一步（转化）：}]
			\[
				\forall x,\ f(x)\leq 5 \iff \max_{x\in[-1,1]}f(x)\leq 5.
			\]
			\par\textbf{理由：} 全称小于转化。
		\item[\textbf{第二步（求$M$）：}]
			对称轴 $x=-\dfrac{a}{2}$，最大值
			\[
				M=
				\begin{cases}
					a+2, & a\ge0\\ 2-a, & a<0
			\end{cases}
				.
			\]
			\par\textbf{理由：} 二次函数开口向上，最大值在端点。
		\item[\textbf{第三步（解不等式）：}]
			$M\leq5$ 得 $
			\begin{aligned}
				&a\ge0,\ a+2\leq5\Rightarrow 0\le a\le3\\ &a<0,\ 2-a\leq5\Rightarrow -3\le a<0
		\end{aligned}
			$，故 $a\in[-3,3]$。
			\par\textbf{理由：} 分段解再合并。
	\end{description}
	\textbf{\textcolor{CExample}{关键结论：}}
	\textbf{$a\in[-3,3]$}
\end{examplebox}
